\documentclass[10pt]{article}
\usepackage[margin=1in]{geometry}
\usepackage{amsmath,amssymb,amsthm,microtype,xcolor}
\usepackage[colorlinks=true,urlcolor=blue!50!black]{hyperref}
\newtheorem{theorem}{Theorem}
\setlength{\emergencystretch}{3em}
\title{An Endpoint-Swap Obstruction for Natural Collatz Cycles}
\author{Collatz proof project}
\date{September 5, 2026}
\begin{document}
\maketitle
\begin{abstract}
We formalize an elementary cancellation obstruction for actual natural
Collatz returns. If the binary words $1u0$ and $0u1$ are realized by returning
natural seeds, then $u$ is empty and the seeds are respectively 1 and 2.
No palindrome, coprimality, or common-orbit assumption is required.
The argument adapts the cancellation underlying Knight's high-cycle
exclusion to the main project's word and orbit definitions. It is not a
proof of Collatz or a claim of mathematical novelty. The word-structure
bridge needed to apply it to all rational mechanical itineraries remains open
in this main-project formalization.
\end{abstract}

\section{Actual return words}
Use the shortcut map $T(n)=n/2$ for even $n$ and $(3n+1)/2$ for odd $n$.
A word $w$ is realized by $n$ when its bits give the parities of the first
$|w|$ orbit values. Let $j(w)$ count its ones and let $C(w)$ denote its
affine correction, defined recursively by
\[
 C(\epsilon)=0,\quad C(0w)=2C(w),\quad
 C(1w)=3^{j(w)}+2C(w).
\]
The verified affine identity is
$2^{|w|}T^{|w|}(n)=3^{j(w)}n+C(w)$.

\begin{theorem}
Let $u$ be any finite binary word. If $x$ realizes $1u0$, $y$ realizes
$0u1$, and both seeds return after $|u|+2$ shortcut steps, then
\[
 u=\epsilon,\qquad x=1,\qquad y=2.
\]
\end{theorem}
The first letter of $1u0$ already forces $x>0$. The two returns need not
belong to the same orbit. For one cyclic orbit, the theorem applies whenever
a shift realizes the second word; a return at one state implies the same
return length at every shifted state.

\section{Cancellation proof}
Set $\ell=|u|$, $j=j(u)+1$, $A=2^{\ell+2}$, and $B=3^j$.
The positive-cycle inequality gives $B<A$, so $D=A-B$ is a positive
odd integer. Write $C_L=C(1u0)$ and $C_R=C(0u1)$.
Direct expansion, valid for every $u$, yields
\[
3C_L+A=C_R+B+2^{\ell+1}.
\]
The two return identities give $C_L=Dx$ and $C_R=Dy$. Substitution gives
\[
D(3x+1)=Dy+2^{\ell+1},\qquad D\mid2^{\ell+1}.
\]
An odd positive divisor of a power of two is 1. Thus $A=B+1$.
If $\ell>0$, then $A$ is divisible by 8, whereas $3^j$ is congruent to
1 or 3 modulo 8. This contradicts $A=B+1$. Consequently $\ell=0$.
The two affine equations now reduce to $4x=3x+1$ and $4y=3y+2$,
giving $x=1$ and $y=2$.

\section{Lean statements and validation}
The new module \texttt{Collatz/EndpointCycles.lean} provides
\begin{itemize}
\item \texttt{endpoint\_cancellation}, the arbitrary-word identity;
\item \texttt{endpoint\_returns\_trivial}, the two-seed theorem;
\item \texttt{endpoint\_cycle\_shift\_trivial}, its single-orbit corollary.
\end{itemize}
All three are included in the selected axiom audit. The module builds with
only standard foundational axiom dependencies; the conditional finite-range
verification in \texttt{Cycles.lean} is not used. The audit is a Lean build
and dependency inspection, not an independent kernel replay.

The exact Python controls in \texttt{analysis/test\_endpoint\_cycles.py}
check the cancellation and paired-integrality conclusion for every middle
of length at most 12. They also find endpoint-swapped rotations for every
reduced rational mechanical slope $p/q$ with $0<p<q\le80$.
These are finite controls, not a proof of the universal rotation statement.

\section{Provenance and remaining bridge}
The cancellation was located in the nested repository
\texttt{tcosmo\_\_Knight2026\_lean}. The relevant file is
\texttt{Knight2026/NoHighCycles.lean}. Its high-cycle proof also uses
Christoffel word structure and rotation preservation of integrality.
The present proof uses the same mathematical idea, with fresh natural-number
arithmetic in the main project's definitions, and imports no nested module.
The relevant published source is Kevin Knight,
\emph{Collatz high cycles do not exist}, Discrete Mathematics
349(3) (2026), article 114812,
\href{https://doi.org/10.1016/j.disc.2025.114812}{doi:10.1016/j.disc.2025.114812}.
The nested summary's article number 114425 and HAL identifier are inaccurate;
the publisher lists 114812 and the challenge catalog links HAL-04261183.

To finish the rational mechanical route, one must prove that a primitive
rational mechanical period has both endpoint-swapped words among its
rotations, and transport those rotations to actual returning states.
Handling nonprimitive periods and eventual tails is also required.
The irrational mechanical exclusion already proved in the main project does
not supply these steps.

The witness property is not universal for arbitrary parity words. For example,
$11100$ has no endpoint-swapped pair among its rotations. Its formal cyclic
value is $19/5$, so it is a rational example illustrating the structural gap,
not an integer Collatz counterexample. Neither general nontrivial cycles nor
unbounded orbits are excluded by this result.
\end{document}
